\documentclass[10pt,a4paper]{amsart}
\usepackage{amsmath,amssymb,amsthm,graphicx}
\usepackage{verbatim}
\newenvironment{reference}{\comment}{\endcomment}
%\newenvironment{reference}{}{}
\newenvironment{slogan}{\comment}{\endcomment}
\newenvironment{history}{\comment}{\endcomment}

% For commutative diagrams we use Xy-pic
\usepackage[all]{xy}

% We use 2cell for 2-commutative diagrams.
\xyoption{2cell}
% Theorem environments.
%
\theoremstyle{plain}
\newtheorem{theorem}[subsection]{Theorem}
\newtheorem{proposition}[subsection]{Proposition}
\newtheorem{lemma}[subsection]{Lemma}

\theoremstyle{definition}
\newtheorem{definition}[subsection]{Definition}
\newtheorem{example}[subsection]{Example}
\newtheorem{exercise}[subsection]{Exercise}
\newtheorem{situation}[subsection]{Situation}

\theoremstyle{remark}
\newtheorem{remark}[subsection]{Remark}
\newtheorem{remarks}[subsection]{Remarks}

\numberwithin{equation}{subsection}

% Macros
%
\def\lim{\mathop{\mathrm{lim}}\nolimits}
\def\colim{\mathop{\mathrm{colim}}\nolimits}
\def\Spec{\mathop{\mathrm{Spec}}}
\def\Hom{\mathop{\mathrm{Hom}}\nolimits}
\def\Ext{\mathop{\mathrm{Ext}}\nolimits}
\def\SheafHom{\mathop{\mathcal{H}\!\mathit{om}}\nolimits}
\def\SheafExt{\mathop{\mathcal{E}\!\mathit{xt}}\nolimits}
\def\Sch{\mathit{Sch}}
\def\Mor{\mathop{\mathrm{Mor}}\nolimits}
\def\Ob{\mathop{\mathrm{Ob}}\nolimits}
\def\Sh{\mathop{\mathit{Sh}}\nolimits}
\def\NL{\mathop{N\!L}\nolimits}
\def\CH{\mathop{\mathrm{CH}}\nolimits}
\def\proetale{{pro\text{-}\acute{e}tale}}
\def\etale{{\acute{e}tale}}
\def\QCoh{\mathit{QCoh}}
\def\Ker{\mathop{\mathrm{Ker}}}
\def\Im{\mathop{\mathrm{Im}}}
\def\Coker{\mathop{\mathrm{Coker}}}
\def\Coim{\mathop{\mathrm{Coim}}}

% Boxtimes
%
\DeclareMathSymbol{\boxtimes}{\mathbin}{AMSa}{"02}

%
% Macros for moduli stacks/spaces
%
\def\QCohstack{\mathcal{QC}\!\mathit{oh}}
\def\Cohstack{\mathcal{C}\!\mathit{oh}}
\def\Spacesstack{\mathcal{S}\!\mathit{paces}}
\def\Quotfunctor{\mathrm{Quot}}
\def\Hilbfunctor{\mathrm{Hilb}}
\def\Curvesstack{\mathcal{C}\!\mathit{urves}}
\def\Polarizedstack{\mathcal{P}\!\mathit{olarized}}
\def\Complexesstack{\mathcal{C}\!\mathit{omplexes}}
% \Pic is the operator that assigns to X its picard group, usage \Pic(X)
% \Picardstack_{X/B} denotes the Picard stack of X over B
% \Picardfunctor_{X/B} denotes the Picard functor of X over B
\def\Pic{\mathop{\mathrm{Pic}}\nolimits}
\def\Picardstack{\mathcal{P}\!\mathit{ic}}
\def\Picardfunctor{\mathrm{Pic}}
\def\Deformationcategory{\mathcal{D}\!\mathit{ef}}
\usepackage{cleveref}
\crefname{originallemma}{original lemma}{original lemmas}
\Crefname{originallemma}{Original lemma}{Original lemmas}
\makeatletter
\newlabel{lemma-covering-gives-surjection}{{0934}{source}{}{sourcepointer-1}{}}
\newlabel{lemma-covering-gives-surjection@cref}{{[originallemma][0][]0934}{[1][0][]source}{}{}{}}
\makeatother
\crefname{originalreference}{original reference}{original references}
\Crefname{originalreference}{Original reference}{Original references}
\makeatletter
\newlabel{sites-lemma-conclude-quasi-compact}{{0D05}{source}{}{sourcepointer-2}{}}
\newlabel{sites-lemma-conclude-quasi-compact@cref}{{[originalreference][0][]0D05}{[1][0][]source}{}{}{}}
\makeatother
\crefname{originalreference}{original reference}{original references}
\Crefname{originalreference}{Original reference}{Original references}
\makeatletter
\newlabel{sites-definition-morphism-coverings}{{00VT}{source}{}{sourcepointer-3}{}}
\newlabel{sites-definition-morphism-coverings@cref}{{[originalreference][0][]00VT}{[1][0][]source}{}{}{}}
\makeatother
\crefname{originalreference}{original reference}{original references}
\Crefname{originalreference}{Original reference}{Original references}
\makeatletter
\newlabel{homology-lemma-check-exactness}{{05AA}{source}{}{sourcepointer-4}{}}
\newlabel{homology-lemma-check-exactness@cref}{{[originalreference][0][]05AA}{[1][0][]source}{}{}{}}
\makeatother
\crefname{originalequation}{original equation}{original equations}
\Crefname{originalequation}{Original equation}{Original equations}
\makeatletter
\newlabel{equation-map-lower-shriek-OU-into-module}{{0G1V}{source}{}{sourcepointer-5}{}}
\newlabel{equation-map-lower-shriek-OU-into-module@cref}{{[originalequation][0][]0G1V}{[1][0][]source}{}{}{}}
\makeatother
\crefname{originalreference}{original reference}{original references}
\Crefname{originalreference}{Original reference}{Original references}
\makeatletter
\newlabel{sites-lemma-compare-sheaf-condition}{{0G1L}{source}{}{sourcepointer-6}{}}
\newlabel{sites-lemma-compare-sheaf-condition@cref}{{[originalreference][0][]0G1L}{[1][0][]source}{}{}{}}
\makeatother
\title{A covering of a quasi-compact object contains a finite subfamily presenting the object by extension by zero}
\author{The Stacks Project Authors}
\date{stacks-project:sites-modules.tex:local-snapshot-20261012; GNU FDL 1.2 or later}
\begin{document}
\maketitle
\noindent\textit{Editorial note.} All delivered mathematical text is the author's own, copied verbatim from the pinned Stacks Project chapter file sites-modules.tex; no statement or proof step was rewritten, reordered or removed.\par\medskip
\noindent\textit{Editorial note.} The author invites the reader to skip the general case; the general argument is nevertheless present in the span and constitutes the counted proof, so the invitation was kept verbatim as part of the author's text.\par\medskip
\noindent\textit{Editorial note.} The counted span cites six external items (lemmas, definitions, an equation), recorded in assembly.external\_references with canonical tags as navigation pointers.\par\medskip
\noindent\textit{Allowed standard prerequisites (editorial).}\begin{itemize}
\item Quasi-compactness conditions for objects of a site and morphisms of coverings. Stacks Project, Sites, Lemma sites-lemma-conclude-quasi-compact (tag 09B1) and Definition sites-definition-morphism-coverings (tag 00VW). Supplies the refinement of a covering with finite index image used in the general case.
\item Criterion for exactness by testing with Hom into sheaves of modules. Stacks Project, Homology, Lemma homology-lemma-check-exactness (tag 010B). Reduces the exactness of the displayed sequence to the corresponding statement after applying Hom.
\end{itemize}\medskip
\ifdefined\hypertarget\hypertarget{sourcepointer-1}{}\fi
\noindent\textit{Source reference.} Original lemma 0934: sources/stacks-chow/tags.txt, source lines 11502--11502. Cross-chapter standard prerequisite cited by the author inside the counted span; it is used as a pointer and its own proof is not part of this item.\par\medskip
\ifdefined\hypertarget\hypertarget{sourcepointer-2}{}\fi
\noindent\textit{Source reference.} Original reference 0D05: sources/stacks-chow/tags.txt, source lines 16300--16300. Cross-chapter standard prerequisite cited by the author inside the counted span; it is used as a pointer and its own proof is not part of this item.\par\medskip
\ifdefined\hypertarget\hypertarget{sourcepointer-3}{}\fi
\noindent\textit{Source reference.} Original reference 00VT: sources/stacks-chow/tags.txt, source lines 1159--1159. Cross-chapter standard prerequisite cited by the author inside the counted span; it is used as a pointer and its own proof is not part of this item.\par\medskip
\ifdefined\hypertarget\hypertarget{sourcepointer-4}{}\fi
\noindent\textit{Source reference.} Original reference 05AA: sources/stacks-chow/tags.txt, source lines 6843--6843. Cross-chapter standard prerequisite cited by the author inside the counted span; it is used as a pointer and its own proof is not part of this item.\par\medskip
\ifdefined\hypertarget\hypertarget{sourcepointer-5}{}\fi
\noindent\textit{Source reference.} Original equation 0G1V: sources/stacks-chow/tags.txt, source lines 20035--20035. Cross-chapter standard prerequisite cited by the author inside the counted span; it is used as a pointer and its own proof is not part of this item.\par\medskip
\ifdefined\hypertarget\hypertarget{sourcepointer-6}{}\fi
\noindent\textit{Source reference.} Original reference 0G1L: sources/stacks-chow/tags.txt, source lines 20026--20026. Cross-chapter standard prerequisite cited by the author inside the counted span; it is used as a pointer and its own proof is not part of this item.\par\medskip
\begin{lemma}
\label{lemma-silly-quasi-compact}
Let $(\mathcal{C}, \mathcal{O})$ be a ringed site. Let
$\mathcal{U} = \{U_i \to U\}_{i \in I}$ be covering of $\mathcal{C}$.
If $U$ is quasi-compact, then there exist a finite subset
$I' \subset I$ such that the sequence
$$
\bigoplus\nolimits_{i, i' \in I'}
j_{U_i \times_U U_{i'}!}\mathcal{O}_{U_i \times_U U_{i'}} \to
\bigoplus\nolimits_{i \in I'}
j_{U_i!}\mathcal{O}_{U_i} \to
j_!\mathcal{O}_U \to 0
$$
is exact.
\end{lemma}

\begin{proof}
This lemma is immediate from Lemma \ref{lemma-covering-gives-surjection}
if $U$ satisfies condition (3) of
Sites, Lemma \ref{sites-lemma-conclude-quasi-compact}.
We urge the reader to skip the proof in the general case.
By definition there exists a covering
$\mathcal{V} = \{V_j \to U\}_{j \in J}$ and a morphism
$\mathcal{V} \to \mathcal{U}$ of families of maps with fixed target
given by $\text{id} : U \to U$, $\alpha : J \to I$, and
$f_j : V_j \to U_{\alpha(j)}$
(see Sites, Definition \ref{sites-definition-morphism-coverings})
such that the image $I' \subset I$ of $\alpha$ is finite.
By Homology, Lemma \ref{homology-lemma-check-exactness}
it suffices to show that for any sheaf of $\mathcal{O}$-modules $\mathcal{F}$
the functor $\Hom_\mathcal{O}(-, \mathcal{F})$ turns the sequence
of the lemma into an exact sequence. By
(\ref{equation-map-lower-shriek-OU-into-module})
we obtain the usual sequence
$$
0 \to
\mathcal{F}(U) \to
\prod\nolimits_{i \in I'} \mathcal{F}(U_i) \to
\prod\nolimits_{i, i' \in I'} \mathcal{F}(U_i \times_U U_{i'})
$$
This is an exact sequence by
Sites, Lemma \ref{sites-lemma-compare-sheaf-condition}
applied to the family of maps $\{U_i \to U\}_{i \in I'}$
which is refined by the covering $\mathcal{V}$.
\end{proof}

\end{document}
